\documentclass[11pt,oneside]{article}

\usepackage[a4paper,margin=27mm,headheight=15pt]{geometry}
\usepackage[T1]{fontenc}
\usepackage[utf8]{inputenc}
\IfFileExists{libertinus.sty}{\usepackage{libertinus}}{\usepackage{lmodern}}
\usepackage{microtype}
\usepackage{amsmath,amssymb,amsthm,mathtools,mathrsfs}
\usepackage{bm}
\usepackage{booktabs,longtable,array,tabularx}
\usepackage{graphicx}
\usepackage{pgfplots}
\usepackage{xcolor}
\usepackage{enumitem}
\usepackage{fancyhdr}
\usepackage{titlesec}
\usepackage{listings}
\usepackage{xurl}
\usepackage{hyperref}
\usepackage[nameinlink,noabbrev]{cleveref}

\definecolor{linkblue}{RGB}{25,75,135}
\definecolor{shadegray}{RGB}{246,247,249}
\definecolor{rulegray}{RGB}{195,201,208}
\pgfplotsset{compat=1.18}
\hypersetup{
  colorlinks=true,
  linkcolor=linkblue,
  citecolor=linkblue,
  urlcolor=linkblue,
  pdftitle={From Exact Dyadic Arithmetic to Endpoint Asymptotics of the Fabius Function},
  pdfauthor={Complementary report to the ProveIt formal development}
}

\pagestyle{fancy}
\fancyhf{}
\fancyhead[L]{\small The Fabius function and Rvachev's up-function}
\fancyhead[R]{\small\nouppercase{\leftmark}}
\fancyfoot[C]{\thepage}
\renewcommand{\headrulewidth}{0.4pt}

\titleformat{\section}{\Large\bfseries}{\thesection}{0.7em}{}
\titleformat{\subsection}{\large\bfseries}{\thesubsection}{0.7em}{}
\titleformat{\subsubsection}{\normalsize\bfseries}{\thesubsubsection}{0.7em}{}
\setcounter{secnumdepth}{3}
\setcounter{tocdepth}{2}
\setlist[itemize]{leftmargin=2em,itemsep=0.25em,topsep=0.4em}
\setlist[enumerate]{leftmargin=2.3em,itemsep=0.25em,topsep=0.4em}
\setlength{\emergencystretch}{2em}

\newtheorem{theorem}{Theorem}[section]
\newtheorem{proposition}[theorem]{Proposition}
\newtheorem{lemma}[theorem]{Lemma}
\newtheorem{corollary}[theorem]{Corollary}
\theoremstyle{definition}
\newtheorem{definition}[theorem]{Definition}
\newtheorem{example}[theorem]{Example}
\theoremstyle{remark}
\newtheorem{remark}[theorem]{Remark}
\newtheorem{warning}[theorem]{Warning}
\newenvironment{proofidea}{\par\noindent\textit{Proof idea.}\ }{\hfill$\square$\par}

\newcommand{\R}{\mathbb R}
\newcommand{\C}{\mathbb C}
\newcommand{\Q}{\mathbb Q}
\newcommand{\N}{\mathbb N}
\newcommand{\Z}{\mathbb Z}
\newcommand{\Up}{\operatorname{up}}
\newcommand{\wt}{\operatorname{wt}}
\newcommand{\e}{\mathrm e}
\newcommand{\ii}{\mathrm i}
\newcommand{\dd}{\,\mathrm d}
\newcommand{\E}{\mathbb E}
\newcommand{\bigO}{\mathrm O}
\newcommand{\smallo}{\mathrm o}
\newcommand{\supp}{\operatorname{supp}}
\newcommand{\qbinom}[3]{\genfrac{[}{]}{0pt}{}{#1}{#2}_{#3}}
\newcommand{\Mmain}{\mathcal M}
\newcommand{\CoeffMass}{\mathcal C}
\newcommand{\Prefix}{\mathcal S}
\newcommand{\TM}{\varepsilon}
\newcommand{\Log}{\operatorname{Log}}

\title{\bfseries From Exact Dyadic Arithmetic to Endpoint Asymptotics of the Fabius Function\\[0.35em]
\large Coefficient extraction, $q$-binomial collapse, phase locking, and the all-orders saddle expansion}
\author{A complementary report to the ProveIt formal development}
\date{Repository snapshot: 25 August 2026}

\begin{document}
\maketitle

\begin{abstract}
The exact arithmetic theory of the Fabius function and its sharp endpoint asymptotic are usually presented as nearly disjoint subjects.  The former gives rational formulae for values at dyadic arguments, including a half-moment recurrence, a finite sum over ordered compositions, a closed Thue--Morse formula, and a two-level formula involving $q$-Pochhammer symbols and Gaussian binomial coefficients.  The latter starts from the negative Laplace product, applies Mellin analysis, introduces the lower Lambert coordinate, and performs a saddle-point inversion.  This report closes the structural gap between them.

The central observation is that the exact reciprocal-power identity
\[
 F(2^{-n})=2^{-\binom n2}[z^n]G(z)
\]
turns the endpoint problem into a coefficient problem for the entire moment function
\[
 G(z)=\prod_{j\ge1}\frac{\e^{z/2^j}-1}{z/2^j}.
\]
If $\lambda_n$ is the large solution of $\lambda_n2^{-\lambda_n}=2^{-n}$, then
$n=\lambda_n-\log_2\lambda_n$.  Vertical coefficient extraction at $z=\lambda_n$ produces a local phase that is, up to an exponentially small tail, \emph{identical} to the local phase in the continuous Bromwich proof.  The apparently different periodic arguments also agree because
\[
 \log_2\lambda_n=\lambda_n-n\equiv\lambda_n\pmod 1.
\]
Consequently the exact dyadic formula re-derives not only the quadratic logarithmic decay but the corrected Gamma--zeta periodic fluctuation, the constant term, the first inverse-Lambert correction, and the complete all-orders expansion.

The $q$-binomial formula leads to the same object after an exact finite cancellation: its Mersenne/$q$-Pochhammer normalizer disappears, and the remaining numerator is a coefficient of $B_m(z)G(z)$, where $B_m$ is a finite Thue--Morse prefix exponential.  This gives the same all-orders asymptotic on every fixed-numerator dyadic ray $m2^{-n}$.  A rescaled Bromwich identity supplies a second, exact explanation: $n$ applications of the refinement equation map the continuous saddle $2^{\lambda_n}$ to the coefficient saddle $\lambda_n$.

The report also separates what follows for arbitrary real $x$ from what is special to exact dyadic rays.  Monotonicity and reciprocal-power values recover the global leading law
\[
 \log F(x)\sim-\frac{(\log x)^2}{2\log2},
\]
but the full $O(1/|\log x|)$ expansion for every real argument requires a uniform two-parameter analysis of the prefix factor $B_m$ when its numerator grows.  Thus the exact formulae do contain the endpoint asymptotic, but they reveal it only after the correct algebraic resummation and saddle scaling.
\end{abstract}

\begingroup
\footnotesize
\tableofcontents
\endgroup
\newpage

\section{The missing bridge}
\label{sec:introduction}

The formal development in the repository \cite{proveit} contains two powerful bodies of results.
On the arithmetic side, it proves exact rational formulae for $F(m/2^n)$, develops half moments and their recurrences, gives a finite composition formula for $F(2^{-n})$, and derives a global Gaussian-binomial formula at base $q=1/2$.  On the analytic side, it derives a corrected sharp expansion of $F(x)$ as $x\downarrow0$, including a nonconstant periodic function with explicit Gamma--zeta Fourier coefficients and an all-orders saddle series.  The accompanying exposition \cite{masterdoc} develops both sides in detail, but it does not turn the exact dyadic identities themselves into a derivation of the endpoint asymptotic.

That omission is conspicuous for three reasons.

\begin{enumerate}
\item The reciprocal powers $2^{-n}$ are simultaneously the cleanest exact arguments and a canonical sequence tending to the endpoint.
\item The exact formulae already contain geometric scales $1,2,4,\ldots$ through Mersenne factors, Thue--Morse blocks, and $q$-binomial coefficients at $q=1/2$; the periodic term in the asymptotic also comes from a geometric dyadic harmonic sum.
\item The exact values provide rational data of arbitrarily high precision, so one expects more than a numerical comparison with an independently derived asymptotic.
\end{enumerate}

The purpose of this report is to exhibit the missing mechanism.  Its shortest form is the chain
\begin{equation}
\boxed{
\begin{aligned}
\text{exact dyadic value}
&\longrightarrow 2^{-\binom n2}[z^n]G(z)\\
&\longrightarrow \text{vertical coefficient saddle at }z=\lambda_n\\
&\longrightarrow \text{negative-product decomposition at }\lambda_n\\
&\longrightarrow \text{the same periodic Lambert asymptotic.}
\end{aligned}}
\label{eq:bridge-chain}
\end{equation}
The algebraic formulae are not merely compatible with the endpoint result.  After resummation they generate the same analytic phase.

\subsection{Main conclusions}

Let
\begin{equation}
 L=\log2,
 \qquad
 \lambda(x)=-\frac1L W_{-1}(-Lx),
 \qquad
 \lambda(x)2^{-\lambda(x)}=x,
\label{eq:lambda-def}
\end{equation}
where $W_{-1}$ is the lower real branch of Lambert's $W$-function \cite{corless1996}.  Define
\begin{equation}
 \Mmain(x)=
 -\frac L2\lambda(x)^2+
 \left(1+\frac L2\right)\lambda(x)
 -\frac12\log\lambda(x)+C_F+\Psi(\lambda(x)),
\label{eq:main-M}
\end{equation}
where $C_F$ and the smooth mean-zero $1$-periodic function $\Psi$ are specified in
\cref{sec:periodic-decomposition}.

The report proves the following statements from the exact dyadic coefficient formula.

\begin{theorem}[Reciprocal-power asymptotic from exact coefficients]
\label{thm:main-inverse}
Let $\lambda_n=\lambda(2^{-n})$, equivalently
\begin{equation}
 \lambda_n-\log_2\lambda_n=n.
\label{eq:lambda-n-equation}
\end{equation}
Then, as $n\to\infty$,
\begin{equation}
 \log F(2^{-n})=\Mmain(2^{-n})+\bigO(\lambda_n^{-1}).
\label{eq:main-inverse-leading}
\end{equation}
More precisely, for every fixed $N\ge1$ there are smooth $1$-periodic functions
$A_j$, independent of the dyadic sequence, such that
\begin{equation}
 \log F(2^{-n})=
 \Mmain(2^{-n})+
 \sum_{j=1}^{N-1}\frac{A_j(\lambda_n)}{\lambda_n^j}
 +\bigO_N(\lambda_n^{-N}).
\label{eq:main-inverse-all}
\end{equation}
The first coefficient is
\begin{equation}
 A_1(u)=\frac1{24}+\frac1{2L}
 -\frac{\Psi''(u)+\Psi'(u)^2}{2L^2}.
\label{eq:A1}
\end{equation}
These are exactly the correction functions in the continuous endpoint expansion.
\end{theorem}

\begin{theorem}[Fixed-numerator dyadic rays]
\label{thm:main-fixed-m}
Fix an integer $m\ge1$.  Put $x_{m,n}=m2^{-n}$ and
$\lambda_{m,n}=\lambda(x_{m,n})$.  Then for every fixed $N\ge1$,
\begin{equation}
 \log F(m2^{-n})=
 \Mmain(m2^{-n})+
 \sum_{j=1}^{N-1}\frac{A_j(\lambda_{m,n})}{\lambda_{m,n}^j}
 +\bigO_{m,N}(\lambda_{m,n}^{-N})
\label{eq:main-fixed-m}
\end{equation}
as $n\to\infty$, with the same functions $A_j$ as in
\cref{thm:main-inverse}.
\end{theorem}

\begin{corollary}[What exact dyadic values imply globally]
\label{cor:global-leading}
The reciprocal-power asymptotic, together with monotonicity of $F$, implies
\begin{equation}
 \frac{\log F(x)}{(\log x)^2}\longrightarrow-\frac1{2\log2}
 \qquad(x\downarrow0).
\label{eq:global-quadratic-law}
\end{equation}
\end{corollary}

The first two statements retain the constant, the periodic fluctuation, and every inverse-power correction.  The third statement is weaker because the mesh between consecutive reciprocal powers is too coarse for a sharp additive asymptotic in $\log F$; this distinction is analyzed in \cref{sec:dedydadicization}.

\section{Exact arithmetic in a coefficient-ready normalization}
\label{sec:exact-layer}

\subsection{The probability law and its entire moment function}

Let $F:[0,1]\to[0,1]$ be the Fabius function in its cumulative-distribution normalization, and let $X$ be a random variable with distribution function $F$.  Equivalently, one may take the convergent random series
\begin{equation}
 X=\sum_{j\ge0}\frac{U_j}{2^{j+1}},
\label{eq:random-series}
\end{equation}
where the $U_j$ are independent uniform random variables on $[0,1]$.  The corresponding density is the translated and rescaled Rvachev up-function.  The support is $[0,1]$, and the law is symmetric about $1/2$.

Define the half moments
\begin{equation}
 d_0=1,
 \qquad
 d_n=n\int_0^1t^{n-1}\Up(t)\dd t=\E X^n
 \quad(n\ge1),
\label{eq:d-moments}
\end{equation}
and the entire moment-generating function
\begin{equation}
 G(z)=\E\e^{zX}=\sum_{n\ge0}d_n\frac{z^n}{n!}.
\label{eq:G-def}
\end{equation}
The refinement equation gives
\begin{equation}
 G(2z)=\frac{\e^z-1}{z}G(z).
\label{eq:G-functional}
\end{equation}
Symmetry of $X$ about $1/2$ gives the second exact relation
\begin{equation}
 G(z)=\e^zG(-z).
\label{eq:G-symmetry}
\end{equation}

It is convenient to remove the factorial from the moments:
\begin{equation}
 a_n=\frac{d_n}{n!}=[z^n]G(z).
\label{eq:a-def}
\end{equation}
The half-moment recurrence becomes
\begin{equation}
 (2^n-1)a_n=
 \sum_{j=0}^{n-1}\frac{a_j}{(n-j+1)!},
 \qquad a_0=1.
\label{eq:a-recurrence}
\end{equation}

\subsection{Reciprocal powers and the triangular factor}

The exact bridge between moments and endpoint values is
\begin{equation}
 \boxed{
 F(2^{-n})=\frac{d_n}{n!\,2^{\binom n2}}
 =2^{-\binom n2}a_n.}
\label{eq:inverse-exact}
\end{equation}
The factor $2^{-\binom n2}$ already accounts for the leading quadratic term
$-(\log2)n^2/2$ in the logarithm.  It does \emph{not}, by itself, account for the
$n\log n$ term, the Lambert phase, the constant, or the periodic fluctuation.  Those are hidden in the much subtler coefficient $a_n$.

Writing $V_n=F(2^{-n})$ and normalizing
$A_n=2^{\binom n2}V_n=a_n$, one obtains the triangular recurrence
\begin{equation}
 A_n=\sum_{k=0}^{n-1}
 \frac{A_k}{(2^n-1)(n-k+1)!}.
\label{eq:normalized-triangular}
\end{equation}
Iterating it over all directed paths from $0$ to $n$ gives the finite composition formula
\begin{equation}
 F(2^{-n})=2^{-\binom n2}
 \sum_{\rho\vDash n}
 \prod_{j=1}^{\operatorname{len}(\rho)}
 \frac1{(2^{s_j}-1)(r_j+1)!},
\label{eq:composition-formula}
\end{equation}
where $\rho=(r_1,\ldots,r_\ell)$ is an ordered composition of $n$ and
$s_j=r_1+\cdots+r_j$.  Formula \eqref{eq:composition-formula} is exact and positive, but it contains $2^{n-1}$ terms.  Its asymptotic structure becomes visible only after the path sum is resummed into a generating function; see \cref{subsec:composition-resummation}.

\subsection{The global $q$-binomial formula}
\label{subsec:qformula}

For
\begin{equation}
 (a;q)_n=\prod_{j=0}^{n-1}(1-aq^j)
\label{eq:qpochhammer}
\end{equation}
and the Gaussian coefficient $\qbinom nkq$, the formal development proves, for the signed global Fabius extension $\mathcal F$,
\begin{align}
 \mathcal F(m/2^n)
 &=\frac{1}{2^{n^2}(1/2;1/2)_n}
 \sum_{k=0}^{n}
 \frac{\qbinom nk{1/2}}
 {4^{\binom k2}(n+k)!}                                    \notag\\
 &\qquad\times
 \sum_{r=0}^{m2^k-1}(-1)^{\wt(r)}
 (r-m2^k+Q)^{n+k},
\label{eq:qbinomial-global}
\end{align}
independently of the complex translation parameter $Q$.  For
$0\le m\le2^n$, the left side is the bounded Fabius value $F(m/2^n)$.

At first sight \eqref{eq:qbinomial-global} seems to offer a direct route to an asymptotic: the prefactor contains $2^{-n^2}$ and the finite product
$(1/2;1/2)_n$.  That reading is misleading.  The double sum undergoes a cancellation of exactly the same scale.  The correct first step is an exact algebraic collapse, not termwise estimation.

Define the Thue--Morse signs and prefix exponential
\begin{equation}
 \TM_h=(-1)^{\wt(h)},
 \qquad
 B_m(z)=\sum_{h=0}^{m-1}\TM_h\e^{(m-1-h)z},
 \qquad B_0=0,
\label{eq:Bm-def}
\end{equation}
and the Mersenne product
\begin{equation}
 \mathfrak m_n=\prod_{\ell=1}^n(2^\ell-1).
\label{eq:mersenne-product}
\end{equation}
The exact finite $q$-binomial extraction theorem gives
\begin{equation}
 \boxed{
 F(m/2^n)=2^{-\binom n2}[z^n]\bigl(B_m(z)G(z)\bigr).}
\label{eq:prefix-coefficient}
\end{equation}
For $m=1$, $B_1=1$ and this reduces to \eqref{eq:inverse-exact}.  The reason the normalizer disappears is
\begin{equation}
 (1/2;1/2)_n=2^{-\binom{n+1}{2}}\mathfrak m_n,
\label{eq:q-mersenne}
\end{equation}
and the centered numerator of \eqref{eq:qbinomial-global} is exactly
\begin{equation}
 \mathfrak m_n[z^n]\bigl(B_m(z)G(z)\bigr).
\label{eq:q-numerator-collapse}
\end{equation}
Thus
\begin{equation}
 \frac{\mathfrak m_n}{2^{n^2}(1/2;1/2)_n}
 =2^{-\binom n2}.
\label{eq:q-normalizer-cancel}
\end{equation}
The $q$-Pochhammer factor supplies the same triangular power already present in
\eqref{eq:inverse-exact}; the remaining endpoint asymptotic resides in a coefficient of the entire product $G$.

\begin{remark}[Why termwise estimates fail]
The individual terms in \eqref{eq:qbinomial-global} are vastly larger than their sum.  The inner Thue--Morse block annihilates low polynomial degrees, and the outer Gaussian row annihilates the remaining lower powers on the geometric node set $-1,-2,\ldots,-2^n$.  Estimating before these exact cancellations discards the quantity one is trying to compute.  Formula \eqref{eq:prefix-coefficient} is therefore not a cosmetic rewrite; it is the asymptotically meaningful normal form.
\end{remark}

\section{Resumming the exact recurrences}
\label{sec:resummation}

\subsection{From the triangular recurrence to the entire product}
\label{subsec:composition-resummation}

Let
\begin{equation}
 \mathscr A(z)=\sum_{n\ge0}a_nz^n=G(z).
\label{eq:A-series}
\end{equation}
Multiplying \eqref{eq:a-recurrence} by $z^n$ and summing over $n\ge1$ gives
\begin{align*}
 \mathscr A(2z)-\mathscr A(z)
 &=\mathscr A(z)\sum_{r\ge1}\frac{z^r}{(r+1)!}\\
 &=\mathscr A(z)\left(\frac{\e^z-1}{z}-1\right),
\end{align*}
so
\begin{equation}
 \mathscr A(2z)=\frac{\e^z-1}{z}\mathscr A(z).
\label{eq:A-functional}
\end{equation}
This is exactly \eqref{eq:G-functional}.  Iteration toward the origin yields the locally uniformly convergent product
\begin{equation}
 \boxed{
 G(z)=\prod_{j=1}^{\infty}
 \frac{\e^{z/2^j}-1}{z/2^j}.}
\label{eq:G-product}
\end{equation}

The path sum \eqref{eq:composition-formula}, the recurrence
\eqref{eq:normalized-triangular}, the functional equation \eqref{eq:A-functional}, and the product \eqref{eq:G-product} are four presentations of one object.  The first two are finite and exact at each $n$; the last two are the forms suited to asymptotics.

\subsection{A probabilistic interpretation of the product}

Each factor is a uniform moment-generating function:
\begin{equation}
 \frac{\e^w-1}{w}=\int_0^1\e^{wu}\dd u.
\label{eq:uniform-mgf}
\end{equation}
Thus \eqref{eq:G-product} is precisely the moment-generating function of the random series \eqref{eq:random-series}.  This gives several useful estimates for free.  For real $r>0$ and real $t$,
\begin{equation}
 |G(r+\ii t)|\le G(r),
\label{eq:G-vertical-bound}
\end{equation}
because $|\e^{(r+\ii t)X}|=\e^{rX}$.  The bound will control the noncentral part of the coefficient integral without a separate zero analysis.

\subsection{Positive and negative arguments}

Define, for $s>0$,
\begin{equation}
 \Lambda(s)=\log G(-s).
\label{eq:Lambda-def}
\end{equation}
Using \eqref{eq:G-symmetry},
\begin{equation}
 \log G(s)=s+\Lambda(s).
\label{eq:positive-negative-log}
\end{equation}
At the product level,
\begin{align}
 G(-s)
 &=\prod_{j=1}^{\infty}
 \frac{1-\e^{-s/2^j}}{s/2^j},                              \label{eq:negative-product}\\
 \Lambda(s)
 &=\sum_{j=1}^{\infty}\kappa(s/2^j),
 \qquad
 \kappa(u)=\log\frac{1-\e^{-u}}u.                         \label{eq:Lambda-harmonic}
\end{align}
This identity is the point where the exact moment product meets the endpoint asymptotic.  No new transform has appeared: the negative Laplace product used in the continuous proof is the negative-real restriction of the entire product forced by the exact recurrence.

\section{The periodic decomposition of the negative product}
\label{sec:periodic-decomposition}

This section records the analytic information about \eqref{eq:Lambda-harmonic} needed by the coefficient saddle.  The derivation is included because it identifies exactly where the $q=1/2$ geometry reappears in the asymptotic.

\subsection{Mellin transform of the dyadic kernel}

For $-1<\Re z<0$,
\begin{equation}
 \widetilde\kappa(z)=\int_0^\infty\kappa(u)u^{z-1}\dd u
 =-\Gamma(z)\zeta(1+z).
\label{eq:kappa-mellin}
\end{equation}
Indeed, $\kappa(u)=-u/2+\bigO(u^2)$ at the origin and
$\kappa(u)=-\log u+\bigO(\e^{-u})$ at infinity.  Integration by parts, followed by the subtracted Bose integral, gives \eqref{eq:kappa-mellin}.

The geometric harmonic sum contributes
\begin{equation}
 \widetilde\Lambda(z)
 =\frac{2^z}{1-2^z}\bigl(-\Gamma(z)\zeta(1+z)\bigr),
 \qquad -1<\Re z<0.
\label{eq:Lambda-mellin}
\end{equation}
The factor
\begin{equation}
 \frac{2^z}{1-2^z}
\label{eq:geometric-mellin-factor}
\end{equation}
is the analytic image of the same dyadic scale that produced the Mersenne factors and the half-base Gaussian coefficients in the exact arithmetic.  Its poles are
\begin{equation}
 \chi_k=\frac{2\pi\ii k}{L},
 \qquad k\in\Z.
\label{eq:chi-k}
\end{equation}
The triple pole at $0$ produces a quadratic polynomial in $\log s$ and a constant; the nonzero poles produce a periodic Fourier series in $\log_2s$.

\subsection{Exact decomposition, mean, and oscillation}

There is a smooth $1$-periodic function $P$ and a positive forward-tail term $E(s)$ such that
\begin{equation}
 \Lambda(s)=
 -\frac{(\log s)^2}{2L}+\frac{\log s}{2}
 +P(\log_2s)+E(s),
\label{eq:Lambda-periodic}
\end{equation}
where
\begin{equation}
 0\le E(s)\le4\e^{-s}
\label{eq:E-bound}
\end{equation}
for all sufficiently large $s$.  Put
\begin{equation}
 \overline P=\int_0^1P(t)\dd t,
 \qquad
 \Psi(t)=P(t)-\overline P.
\label{eq:Psi-def}
\end{equation}
Then
\begin{equation}
 \Psi(t)=\sum_{k\ne0}\widehat\Psi(k)\e^{2\pi\ii kt},
 \qquad
 \widehat\Psi(k)=
 -\frac{\Gamma(-\chi_k)\zeta(1-\chi_k)}L.
\label{eq:Psi-Fourier}
\end{equation}
The Fourier coefficients decay exponentially in $|k|$, so $\Psi$ is real analytic.  It is nonconstant, although its amplitude is only a few parts in $10^6$.

The finite part at the Mellin origin is
\begin{equation}
 A_0=\frac{6\gamma^2+12\gamma_1-\pi^2}{12},
\label{eq:A0}
\end{equation}
where $\gamma$ is Euler's constant and $\gamma_1$ is the first Stieltjes constant.  One has
\begin{equation}
 \overline P=\frac{A_0}{L}-\frac L{12},
\label{eq:Pbar}
\end{equation}
and the endpoint constant after the Gaussian normalization is
\begin{equation}
 C_F=\overline P-\frac12\log(2\pi)
 =\frac{A_0}{L}-\frac{7L}{12}-\frac12\log\pi
 \approx-2.02798389863885942397.
\label{eq:CF}
\end{equation}

\begin{remark}[The same dyadic geometry twice]
In the exact formula, the geometric nodes $-1,-2,\ldots,-2^n$ are isolated by the finite half-base $q$-binomial theorem.  In the asymptotic, the infinite scale family $s/2,s/4,\ldots$ is summed by the Mellin factor $2^z/(1-2^z)$.  Finite $q$-interpolation and Mellin periodicity are therefore two transforms of the same refinement grid.
\end{remark}

\section{Vertical coefficient extraction}
\label{sec:vertical-coeff}

The usual Cauchy coefficient integral on a circle is enough for a saddle analysis.  A vertical representation is better here because it makes the coefficient phase identical to the continuous Bromwich phase.

\begin{lemma}[Vertical extraction of moments]
\label{lem:vertical-extraction}
For every $n\ge1$ and every $c>0$,
\begin{equation}
 a_n=[z^n]G(z)=
 \frac1{2\pi\ii}
 \int_{c-\ii\infty}^{c+\ii\infty}
 \frac{G(z)}{z^{n+1}}\dd z.
\label{eq:vertical-extraction}
\end{equation}
The integral is absolutely convergent.
\end{lemma}

\begin{proof}
On the line $\Re z=c$, \eqref{eq:G-vertical-bound} gives
$|G(z)|\le G(c)$, while $|z|^{-n-1}$ is integrable for $n\ge1$.
Using $G(z)=\E\e^{zX}$, Fubini's theorem and the classical inverse-Laplace formula give
\begin{align*}
 \frac1{2\pi\ii}\int_{c-\ii\infty}^{c+\ii\infty}
 \frac{G(z)}{z^{n+1}}\dd z
 &=\E\left[
 \frac1{2\pi\ii}\int_{c-\ii\infty}^{c+\ii\infty}
 \frac{\e^{zX}}{z^{n+1}}\dd z\right]\\
 &=\E\frac{X^n}{n!}=\frac{d_n}{n!}=a_n.
\end{align*}
\end{proof}

Combining \eqref{eq:inverse-exact} and \eqref{eq:vertical-extraction} gives the exact contour formula
\begin{equation}
 F(2^{-n})=
 \frac{2^{-\binom n2}}{2\pi\ii}
 \int_{c-\ii\infty}^{c+\ii\infty}
 \frac{G(z)}{z^{n+1}}\dd z.
\label{eq:inverse-vertical}
\end{equation}
This is obtained entirely from the exact moment identity; no inversion formula for $F$ has been used.

\section{The reciprocal-power saddle}
\label{sec:inverse-saddle}

\subsection{The Lambert coordinate is forced by coefficient balance}

Let $\lambda=\lambda_n$ be the large solution of \eqref{eq:lambda-n-equation}.  It is not introduced by analogy with the continuous proof.  It is the leading solution of the coefficient saddle equation.

Indeed, by \eqref{eq:positive-negative-log} and \eqref{eq:Lambda-periodic},
\begin{equation}
 \log G(r)=
 r-\frac{(\log r)^2}{2L}+\frac{\log r}{2}
 +P(\log_2r)+E(r).
\label{eq:logG-positive}
\end{equation}
The logarithmic derivative satisfies
\begin{equation}
 r\frac{G'(r)}{G(r)}
 =r-\log_2r+\frac12+\frac{\Psi'(\log_2r)}L+\bigO(r\e^{-r}).
\label{eq:coefficient-saddle-exactish}
\end{equation}
The exact Hayman radius differs by only $\bigO(1)$ from the solution of
$r-\log_2r=n$, and using the latter changes the logarithm of the coefficient only by
$\bigO(1/r)$.  The Lambert coordinate is therefore intrinsic to the coefficient problem.

\subsection{Exact coefficient saddle mass}

Set $c=\lambda$ in \eqref{eq:inverse-vertical} and write
$z=\lambda(1+\ii\theta)$.  Define
\begin{equation}
 \CoeffMass_n=
 \sqrt{\frac{\lambda}{2\pi}}
 \int_{-\infty}^{\infty}
 \frac{G(\lambda(1+\ii\theta))}{G(\lambda)}
 (1+\ii\theta)^{-n-1}\dd\theta.
\label{eq:Cn-def}
\end{equation}
Then
\begin{equation}
 F(2^{-n})=
 \frac{2^{-\binom n2}G(\lambda)\lambda^{-n}}
 {\sqrt{2\pi\lambda}}\CoeffMass_n.
\label{eq:exact-coefficient-saddle}
\end{equation}
Using $G(\lambda)=\exp(\lambda+\Lambda(\lambda))$,
\begin{equation}
 \log F(2^{-n})=
 -L\binom n2+\lambda+\Lambda(\lambda)-n\log\lambda
 -\frac12\log(2\pi\lambda)+\log\CoeffMass_n.
\label{eq:exact-log-coefficient}
\end{equation}

\subsection{Phase locking}
\label{subsec:phase-locking}

The Lambert equation gives the exact integer relation
\begin{equation}
 \log_2\lambda=\lambda-n.
\label{eq:phase-lock}
\end{equation}
Since $\Psi$ is $1$-periodic,
\begin{equation}
 \boxed{\Psi(\log_2\lambda)=\Psi(\lambda).}
\label{eq:Psi-phase-lock}
\end{equation}
This identity is the missing connection between the phase produced by the moment product and the phase used in the endpoint expansion.  A naive coefficient estimate sees the periodic environment at $\log_2\lambda$; the exact dyadic saddle equation shifts that argument by the integer $n$.

Substitute \eqref{eq:Lambda-periodic} into \eqref{eq:exact-log-coefficient}.  Put
$q=\log\lambda=L(\lambda-n)$.  The nonconstant part becomes $\Psi(\lambda)$ by \eqref{eq:Psi-phase-lock}, and the polynomial terms satisfy the exact cancellation
\begin{align}
 &-\frac L2(n^2-n)+\lambda-
 \frac{q^2}{2L}+\frac q2-nq                         \notag\\
 &\hspace{4em}=-\frac L2\lambda^2+
 \left(1+\frac L2\right)\lambda.
\label{eq:key-algebra-cancellation}
\end{align}
Together with \eqref{eq:CF}, this gives an exact separation analogous to the one in the continuous saddle proof.

\begin{proposition}[Exact dyadic separation]
\label{prop:exact-dyadic-separation}
For $\lambda=\lambda_n$,
\begin{equation}
 \boxed{
 \log F(2^{-n})-\Mmain(2^{-n})
 =E(\lambda)+\log\CoeffMass_n.}
\label{eq:exact-dyadic-separation}
\end{equation}
\end{proposition}

All algebraic normalization, the constant, and the periodic term have now been extracted from an exact rational value.  The only remaining task is to show that the dimensionless coefficient saddle mass is $1+\bigO(1/\lambda)$ and to continue its expansion.

\subsection{The local phase is the continuous phase}
\label{subsec:phase-identity}

Let
\begin{equation}
 u=\Log(1+\ii\theta).
\label{eq:u-def}
\end{equation}
On a fixed neighborhood of $\theta=0$, the complex continuation of
\eqref{eq:Lambda-periodic} is valid.  The exponent in \eqref{eq:Cn-def} is
\begin{equation}
 \Phi^{\mathrm{coef}}_\lambda(\theta)=
 \log\frac{G(\lambda(1+\ii\theta))}{G(\lambda)}-(n+1)u.
\label{eq:coef-phase-def}
\end{equation}
Using $G(z)=\e^zG(-z)$, the periodic decomposition, and
$n+\log_2\lambda=\lambda$, one obtains
\begin{align}
 \Phi^{\mathrm{coef}}_\lambda(\theta)
 &=\lambda(\ii\theta-u)-\frac u2-\frac{u^2}{2L}
 +\Psi\!\left(\lambda+\frac uL\right)-\Psi(\lambda)
 +\mathcal E^{\mathrm{coef}}_\lambda(\theta),
\label{eq:coefficient-local-phase}
\end{align}
where $\mathcal E^{\mathrm{coef}}_\lambda(0)=0$ and every derivative of the error is
$\bigO(\e^{-\lambda})$ on a fixed central neighborhood.

Formula \eqref{eq:coefficient-local-phase} is the decisive identity.  Apart from an exponentially small tail, it is exactly the phase obtained by scaling the continuous Bromwich integral at its saddle.  Therefore every local Gaussian polynomial and every all-orders correction coefficient must coincide.

\subsection{Gaussian domination and the leading mass}

The central scale is $\theta=v/\sqrt\lambda$.  Since
\begin{equation}
 \Log(1+\ii\theta)=
 \ii\theta+\frac{\theta^2}{2}-\frac{\ii\theta^3}{3}
 -\frac{\theta^4}{4}+\bigO(\theta^5),
\label{eq:log-expansion}
\end{equation}
the leading term in \eqref{eq:coefficient-local-phase} is
$-\lambda\theta^2/2$.

The tails are especially simple in the vertical representation.  By
\eqref{eq:G-vertical-bound},
\begin{equation}
 \left|
 \frac{G(\lambda(1+\ii\theta))}{G(\lambda)}
 (1+\ii\theta)^{-n-1}
 \right|
 \le(1+\theta^2)^{-(n+1)/2}.
\label{eq:coefficient-tail-bound}
\end{equation}
Since $n\sim\lambda$, this gives Gaussian decay near the origin and exponential decay outside any fixed neighborhood.  Taylor expansion on a window growing more slowly than $\lambda^{1/6}$, followed by \eqref{eq:coefficient-tail-bound}, gives
\begin{equation}
 \CoeffMass_n=1+\bigO(\lambda^{-1}).
\label{eq:Cn-leading}
\end{equation}
The order $\lambda^{-1/2}$ term vanishes after integration because it is odd and purely imaginary.

Combining \eqref{eq:exact-dyadic-separation}, \eqref{eq:E-bound}, and
\eqref{eq:Cn-leading} proves \eqref{eq:main-inverse-leading}.

\subsection{First correction}

Substitute $\theta=v/\sqrt\lambda$ into \eqref{eq:coefficient-local-phase}.  Through order
$\lambda^{-1}$,
\begin{equation}
 \Phi^{\mathrm{coef}}_\lambda(v/\sqrt\lambda)
 =-\frac{v^2}{2}+\lambda^{-1/2}P_1(v;\lambda)
 +\lambda^{-1}P_2(v;\lambda)
 +\bigO\!\left(\lambda^{-3/2}(1+|v|^9)\right),
\label{eq:phase-first-expansion}
\end{equation}
where
\begin{align}
 P_1(v;u)
 &=\ii\left[
 \left(\frac{\Psi'(u)}L-\frac12\right)v+\frac{v^3}{3}
 \right],                                                   \label{eq:P1}\\
 P_2(v;u)
 &=\frac{v^4}{4}+
 \left[-\frac14+\frac1{2L}+\frac{\Psi'(u)}{2L}
 -\frac{\Psi''(u)}{2L^2}\right]v^2.                       \label{eq:P2}
\end{align}
If $Z$ is standard normal, the relative coefficient of $1/\lambda$ is
\begin{equation}
 \E\left[P_2(Z;u)+\frac12P_1(Z;u)^2\right]
 =\frac1{24}+\frac1{2L}
 -\frac{\Psi''(u)+\Psi'(u)^2}{2L^2}.
\label{eq:first-Gaussian-contraction}
\end{equation}
This is precisely $A_1(u)$ in \eqref{eq:A1}.  Since the logarithm of
$1+A_1/\lambda+\bigO(\lambda^{-2})$ has the same first coefficient,
\begin{equation}
 \log\CoeffMass_n=
 \frac{A_1(\lambda_n)}{\lambda_n}+\bigO(\lambda_n^{-2}).
\label{eq:Cn-first}
\end{equation}
Thus the first correction in the continuous asymptotic has been recovered from the exact coefficient formula.

\subsection{All orders}

At order $N$, expand \eqref{eq:coefficient-local-phase} after the central scaling through
$\lambda^{-N}$.  The coefficients are polynomials in $v$ whose scalar coefficients are differential polynomials in
\begin{equation}
 \Psi'(\lambda),\Psi''(\lambda),\ldots,\Psi^{(2N)}(\lambda).
\end{equation}
After exponentiation, all odd powers of $\lambda^{-1/2}$ integrate to zero by parity.  Gaussian moments convert the remaining terms into integer powers of $1/\lambda$.  The tail bound \eqref{eq:coefficient-tail-bound} justifies termwise integration and gives an arbitrarily small remainder.  Consequently
\begin{equation}
 \log\CoeffMass_n=
 \sum_{j=1}^{N-1}\frac{A_j(\lambda_n)}{\lambda_n^j}
 +\bigO_N(\lambda_n^{-N}),
\label{eq:Cn-all-orders}
\end{equation}
with the same $A_j$ as in the continuous saddle expansion, because the local phase is the same.  This completes the proof of \cref{thm:main-inverse}.

\section{The elementary $n$-scale expansion}
\label{sec:n-scale}

The Lambert form is the cleanest because it preserves the exact periodic argument.  Expanding the phase in the integer parameter gives a more familiar formula.  The large solution of
\eqref{eq:lambda-n-equation} satisfies
\begin{equation}
 \lambda_n=n+\frac{\log n}{L}
 +\frac{\log n}{L^2n}
 +\frac{\log n-\tfrac12\log^2n}{L^3n^2}
 +\bigO\!\left(\frac{\log^3n}{n^3}\right).
\label{eq:lambda-n-expansion}
\end{equation}
Substitution into \eqref{eq:main-inverse-leading}, with the exact
$\Psi(\lambda_n)$ retained, gives
\begin{align}
 \log F(2^{-n})
 ={}&-\frac L2n^2-n\log n+
 \left(1+\frac L2\right)n
 -\frac{\log^2n}{2L}+C_F                              \notag\\
 &-\frac{\log^2n}{2L^2n}
 +\Psi(\lambda_n)+\bigO(n^{-1}).
\label{eq:integer-dyadic-expansion}
\end{align}
This is now a consequence of the exact rational coefficient sequence rather than an independent specialization of the continuous theorem.

Formula \eqref{eq:integer-dyadic-expansion} also clarifies the roles of the factors in
\eqref{eq:inverse-exact}:
\begin{itemize}
\item $2^{-\binom n2}$ gives the quadratic term and part of the linear term;
\item $a_n=d_n/n!$ supplies the cancellation that creates $-n\log n$ and the remaining drift;
\item the dyadic product for $G$ supplies $C_F$ and $\Psi$;
\item Gaussian coefficient extraction supplies $-\tfrac12\log\lambda_n$ and the functions $A_j$.
\end{itemize}

\section{Fixed-numerator dyadic rays}
\label{sec:fixed-numerator}

The coefficient form \eqref{eq:prefix-coefficient} makes it possible to repeat the derivation for $m2^{-n}$ with fixed $m$.

\subsection{The prefix factor near the saddle}

Write
\begin{equation}
 B_m(z)=\e^{(m-1)z}\Prefix_m(z),
 \qquad
 \Prefix_m(z)=\sum_{h=0}^{m-1}\TM_h\e^{-hz}.
\label{eq:Sm-def}
\end{equation}
For $\Re z=r>0$,
\begin{equation}
 |\Prefix_m(z)-1|
 \le\sum_{h\ge1}\e^{-hr}
 =\frac{\e^{-r}}{1-\e^{-r}}.
\label{eq:Sm-bound}
\end{equation}
For fixed $m$ and $r\to\infty$, the entire prefix is therefore dominated by its first exponential:
\begin{equation}
 B_m(z)=\e^{(m-1)z}\left(1+\bigO_m(\e^{-r})\right),
\label{eq:Bm-dominant}
\end{equation}
uniformly on vertical lines and, with derivatives, on a fixed central window.

\subsection{The ray saddle and its phase lock}

Let
\begin{equation}
 x=m2^{-n},
 \qquad
 \lambda=\lambda(x),
 \qquad
 r=\frac\lambda m.
\label{eq:fixed-m-parameters}
\end{equation}
The Lambert equation implies
\begin{equation}
 n=\lambda-\log_2\lambda+\log_2m
 =mr-\log_2r.
\label{eq:fixed-m-saddle-relation}
\end{equation}
Equivalently,
\begin{equation}
 \log_2r=\lambda-n,
 \qquad
 \Psi(\log_2r)=\Psi(\lambda).
\label{eq:fixed-m-phase-lock}
\end{equation}
Thus the exact periodic phase locks in the same way as for $m=1$.

For fixed $m$ and all sufficiently large $n$, $\Prefix_m(r)>0$, so the following real logarithms are well defined.  Apply vertical coefficient extraction to
$H_m(z)=B_m(z)G(z)$:
\begin{equation}
 [z^n]H_m(z)=
 \frac1{2\pi\ii}\int_{r-\ii\infty}^{r+\ii\infty}
 \frac{H_m(z)}{z^{n+1}}\dd z.
\label{eq:Hm-vertical}
\end{equation}
Define the normalized ray saddle mass
\begin{equation}
 \CoeffMass_{m,n}=
 \sqrt{\frac\lambda{2\pi}}
 \int_{-\infty}^{\infty}
 \frac{B_m(r(1+\ii\theta))G(r(1+\ii\theta))}
 {B_m(r)G(r)}
 (1+\ii\theta)^{-n-1}\dd\theta.
\label{eq:Cmn-def}
\end{equation}
Then the exact coefficient formula gives
\begin{align}
 \log F(m2^{-n})
 ={}&-L\binom n2+\log B_m(r)+\log G(r)-n\log r    \notag\\
 &-\frac12\log(2\pi\lambda)+\log\CoeffMass_{m,n}.
\label{eq:fixed-m-exact-log}
\end{align}
Using \eqref{eq:Sm-def} and \eqref{eq:positive-negative-log},
\begin{equation}
 \log B_m(r)+\log G(r)
 =\lambda+\Lambda(r)+\log\Prefix_m(r).
\label{eq:BmG-real}
\end{equation}
The same algebra as \eqref{eq:key-algebra-cancellation}, now with
$\log r=L(\lambda-n)$, yields the exact separation
\begin{equation}
 \boxed{
 \log F(m2^{-n})-\Mmain(m2^{-n})
 =E(r)+\log\Prefix_m(r)+\log\CoeffMass_{m,n}.}
\label{eq:fixed-m-exact-separation}
\end{equation}

\subsection{Identity of local phases}

On the central window, \eqref{eq:Bm-dominant} gives
\begin{align*}
 &\log\frac{B_m(r(1+\ii\theta))G(r(1+\ii\theta))}
 {B_m(r)G(r)}-(n+1)u\\
 &\quad=\lambda(\ii\theta-u)-\frac u2-\frac{u^2}{2L}
 +\Psi\!\left(\lambda+\frac uL\right)-\Psi(\lambda)
 +\mathcal E_{m,\lambda}(\theta),
\end{align*}
where every derivative of $\mathcal E_{m,\lambda}$ is
$\bigO_m(\e^{-\lambda/m})$.  This is again
\eqref{eq:coefficient-local-phase}.  The vertical tails are bounded by the denominator
$(1+\theta^2)^{-(n+1)/2}$ times a fixed $m$-dependent constant.  Therefore
\begin{equation}
 \log\CoeffMass_{m,n}=
 \sum_{j=1}^{N-1}\frac{A_j(\lambda)}{\lambda^j}
 +\bigO_{m,N}(\lambda^{-N}).
\label{eq:Cmn-all-orders}
\end{equation}
The two additional terms in \eqref{eq:fixed-m-exact-separation} are exponentially small for fixed $m$.  This proves \cref{thm:main-fixed-m}.

\subsection{A moderate-numerator extension}

The estimate \eqref{eq:Sm-bound} is independent of $m$.  Its derivatives on the central scale introduce only fixed powers of the reduced saddle
$r=\lambda/m$.  Consequently the same argument remains uniform whenever the exponential prefix error dominates every such polynomial loss.

\begin{proposition}[A useful uniform regime]
\label{prop:moderate-m}
Let $m=m(n)$, put $\lambda=\lambda(m2^{-n})$, and suppose
\begin{equation}
 \frac{\lambda/m}{\log\lambda}\longrightarrow\infty.
\label{eq:moderate-m-condition}
\end{equation}
Then, for every fixed $N\ge1$, the expansion in \eqref{eq:main-fixed-m} holds through order $\lambda^{-N}$ uniformly along that sequence.
\end{proposition}

Indeed, every fixed collection of derivatives of the prefix correction is bounded by
$\bigO_N(r^{K_N}\e^{-r})$ for some finite $K_N$, and \eqref{eq:moderate-m-condition} makes this smaller than every power of $1/\lambda$.  The condition is equivalent to
$m=o(\lambda/\log\lambda)$.  It does not cover the fine dyadic meshes needed to approximate every real $x$ at the sharp logarithmic scale; see
\cref{sec:dedydadicization}.

\section{An exact Bromwich-to-coefficient collapse}
\label{sec:bromwich-collapse}

The coefficient derivation above is already sufficient.  A second identity shows directly how the continuous endpoint saddle and the exact dyadic coefficient saddle are the same object viewed before and after $n$ refinement steps.

\subsection{Iterating the negative refinement equation}

From \eqref{eq:G-functional},
\begin{equation}
 G(-2z)=\frac{1-\e^{-z}}zG(-z).
\label{eq:negative-refinement}
\end{equation}
After $n$ iterations,
\begin{equation}
 G(-2^nz)=
 2^{-\binom n2}z^{-n}G(-z)
 \prod_{j=0}^{n-1}(1-\e^{-2^jz}).
\label{eq:negative-refinement-iterated}
\end{equation}
The finite product is the Thue--Morse block polynomial
\begin{equation}
 P_n(z)=\prod_{j=0}^{n-1}(1-\e^{-2^jz})
 =\sum_{h=0}^{2^n-1}\TM_h\e^{-hz}.
\label{eq:Pn-TM}
\end{equation}

For $0<x<1$, the Laplace--Stieltjes inversion formula is
\begin{equation}
 F(x)=\frac1{2\pi\ii}
 \int_{R-\ii\infty}^{R+\ii\infty}
 \frac{\e^{sx}G(-s)}s\dd s,
 \qquad R>0.
\label{eq:Bromwich-F}
\end{equation}
Set $x=2^{-n}$ and substitute $s=2^nz$.  Then
\begin{align}
 F(2^{-n})
 &=\frac1{2\pi\ii}\int
 \frac{\e^zG(-2^nz)}z\dd z                              \notag\\
 &=\frac{2^{-\binom n2}}{2\pi\ii}
 \int
 \frac{G(z)P_n(z)}{z^{n+1}}\dd z,
\label{eq:Bromwich-rescaled}
\end{align}
where \eqref{eq:G-symmetry} was used in the second line.

\subsection{Why the finite product disappears exactly}

Insert \eqref{eq:Pn-TM} and $G(z)=\E\e^{zX}$ into
\eqref{eq:Bromwich-rescaled}.  The inverse-Laplace kernel gives
\begin{align}
 F(2^{-n})
 &=2^{-\binom n2}
 \E\left[
 \sum_{h=0}^{2^n-1}\TM_h\frac{(X-h)_+^n}{n!}
 \right].
\label{eq:positive-part-collapse}
\end{align}
Because $0\le X\le1$, all terms with $h\ge1$ vanish almost surely.  The $h=0$ term is
$\E X^n/n!=a_n$, and \eqref{eq:inverse-exact} follows.  Thus the exact reciprocal-power bridge itself can be read as a support-induced collapse of the rescaled Bromwich integral.

Near the saddle this collapse has an even simpler form.  On the line $\Re z=\lambda_n$,
\begin{equation}
 P_n(z)=1+\bigO(\e^{-\lambda_n})
\label{eq:Pn-near-one}
\end{equation}
uniformly in $\Im z$, because
$\sum_{j\ge0}\e^{-2^j\lambda_n}=\bigO(\e^{-\lambda_n})$.
The continuous saddle radius is
\begin{equation}
 R=\frac{\lambda_n}{2^{-n}}=2^n\lambda_n=2^{\lambda_n},
\label{eq:continuous-vs-coeff-saddle}
\end{equation}
where the last equality is \eqref{eq:lambda-n-equation}.  The substitution
$s=2^nz$ maps $R$ exactly to $z=\lambda_n$, and
\eqref{eq:negative-refinement-iterated} maps the continuous integrand to the coefficient integrand, up to the beyond-all-orders factor $P_n(z)$.  This is the most literal form of the connection sought in this report.

\section{What the different exact formulae contribute}
\label{sec:formula-audit}

The exact formulae are not interchangeable as asymptotic tools.  Their roles can be summarized as follows.

\begingroup\small
\begin{longtable}{@{}
>{\raggedright\arraybackslash}p{0.24\textwidth}
>{\raggedright\arraybackslash}p{0.31\textwidth}
>{\raggedright\arraybackslash}p{0.37\textwidth}@{}}
\toprule
Exact representation & What it exposes directly & What must be done before asymptotics\\
\midrule
\endhead
Half-moment identity
$F(2^{-n})=2^{-\binom n2}d_n/n!$
& Separates the universal triangular power of two and identifies the residual as a moment coefficient.
& Replace $d_n/n!$ by $[z^n]G(z)$ and perform coefficient extraction.\\[0.6em]
Triangular recurrence
\eqref{eq:normalized-triangular}
& Gives exact rational computation and a positive recursion.
& Sum over $n$ to recover $G(2z)=((\e^z-1)/z)G(z)$.\\[0.6em]
Composition/path formula
\eqref{eq:composition-formula}
& Gives a finite nonrecursive sum and a combinatorial interpretation of the triangular factor.
& Resum the $2^{n-1}$ paths into the same functional equation; termwise asymptotics obscure the collective saddle.\\[0.6em]
$q$-binomial formula
\eqref{eq:qbinomial-global}
& Makes the geometric interpolation and two levels of exact cancellation explicit.
& Perform the finite Gaussian-binomial extraction first.  The Mersenne and $q$-Pochhammer factors then cancel, leaving \eqref{eq:prefix-coefficient}.\\[0.6em]
Closed Thue--Morse formula
& Gives every dyadic numerator and reveals binary blocks.
& Convert it to $[z^n](B_mG)$; for fixed $m$, the leading exponential in $B_m$ supplies the ray saddle.\\[0.6em]
Bit/Taylor recursion
& Is efficient for exact evaluation at a specified dyadic argument.
& It is locally subtractive and hides the global saddle; it is best used for numerical verification, not as a termwise asymptotic expansion.\\
\bottomrule
\end{longtable}
\endgroup

\subsection{Why the $q$-Pochhammer factor is not the periodic term}

A tempting but incorrect heuristic is to take logarithms of
$(1/2;1/2)_n$ and seek the endpoint constant or oscillation there.  The finite product converges to a nonzero constant after its explicit triangular normalization, but in
\eqref{eq:qbinomial-global} it is paired with a numerator containing the same Mersenne product.  Their exact cancellation occurs before any limit.  The nonconstant periodic term instead comes from the \emph{infinite} product for $G$ and the nonzero imaginary-axis poles of the Mellin factor \eqref{eq:geometric-mellin-factor}.

\subsection{Why the composition sum must be resummed}

The composition weights are positive, so one might hope for a direct probabilistic limit theorem on random compositions.  The difficulty is that the endpoint index $s_j$ appears in the Mersenne factor $2^{s_j}-1$, while the jump size $r_j$ appears in $(r_j+1)!$.  The dominant path length and jump profile drift with $n$, and no fixed finite family of compositions controls the sum.  The generating function performs the exact renewal resummation and turns that drifting ensemble into the single analytic saddle of $G$.

\section{Exact-rational numerical checks}
\label{sec:numerics}

The following computations use the recurrence \eqref{eq:a-recurrence} in exact rational arithmetic.  Floating-point conversion occurs only after the rational value has been formed.  The periodic function and its derivatives are evaluated from the Fourier series \eqref{eq:Psi-Fourier}.

Let
\begin{equation}
 R_n=\log F(2^{-n})-\Mmain(2^{-n}).
\label{eq:Rn-def}
\end{equation}
The first-order prediction is
$R_n=A_1(\lambda_n)/\lambda_n+\bigO(\lambda_n^{-2})$.

\begin{table}[htbp]
\centering
\small
\begin{tabular}{@{}rrrrrr@{}}
\toprule
$n$ & $\lambda_n$ & $\log F(2^{-n})$ & $R_n$ & $\lambda_nR_n$ &
$\lambda_n^2(R_n-A_1/\lambda_n)$\\
\midrule
12 & 16.0000000000 & $-70.54224031962472$ & 0.04969809560 & 0.7951695295 & 0.5134917380\\
24 & 28.8505257022 & $-253.3041353903491$ & 0.02706676636 & 0.7808904386 & 0.5161104586\\
48 & 53.7481430054 & $-932.8335151659519$ & 0.01437385195 & 0.7725678503 & 0.5168295087\\
96 & 102.682040054 & $-3520.201176273548$ & 0.007479125975 & 0.7679719130 & 0.5174474457\\
\bottomrule
\end{tabular}
\caption{Reciprocal-power values computed exactly.  The scaled residual
$\lambda_nR_n$ approaches the periodic first coefficient $A_1(\lambda_n)$, which is approximately $0.763$ on these phases.  After subtracting that term, multiplication by $\lambda_n^2$ produces a bounded second-order quantity.}
\label{tab:inverse-numerics}
\end{table}

The fixed-numerator theorem can be tested independently by extracting
$[z^n](B_mG)$ exactly.  The table reports
\begin{equation}
 R_{m,n}=\log F(m2^{-n})-\Mmain(m2^{-n}).
\end{equation}

\begin{table}[htbp]
\centering
\small
\begin{tabular}{@{}rrrrrr@{}}
\toprule
$m$ & $n$ & $\lambda_{m,n}$ & $\lambda_{m,n}R_{m,n}$ &
$A_1(\lambda_{m,n})$ &
$\lambda_{m,n}^2(R_{m,n}-A_1/\lambda_{m,n})$\\
\midrule
1 & 20 & 24.62186833 & 0.7836542295 & 0.7629268732 & 0.5104\\
1 & 40 & 45.50804986 & 0.7742689337 & 0.7629490545 & 0.5151\\
1 & 80 & 86.43351899 & 0.7689739455 & 0.7629823190 & 0.5179\\
\addlinespace
3 & 20 & 22.93448405 & 0.7855611616 & 0.7630465477 & 0.5164\\
3 & 40 & 43.87020711 & 0.7748252957 & 0.7630120722 & 0.5182\\
3 & 80 & 84.82139378 & 0.7691025485 & 0.7629858145 & 0.5188\\
\addlinespace
5 & 20 & 22.14711904 & 0.7864070706 & 0.7631005942 & 0.5162\\
5 & 40 & 43.10795409 & 0.7751701258 & 0.7631010621 & 0.5203\\
5 & 80 & 84.07161885 & 0.7692992153 & 0.7630968112 & 0.5214\\
\bottomrule
\end{tabular}
\caption{Exact coefficient checks on the rays $m2^{-n}$.  The same first correction
$A_1$ works for $m=1,3,5$, and the twice-scaled remainder stays bounded, as predicted by
\cref{thm:main-fixed-m}.}
\label{tab:fixed-m-numerics}
\end{table}

The numerical convergence is deliberately slow at first order: if
$R=A_1/\lambda+A_2/\lambda^2+\cdots$, then
$\lambda R=A_1+A_2/\lambda+\cdots$.  The tables show exactly this drift and then stabilize it by subtracting $A_1/\lambda$.

\section{From exact dyadics to arbitrary real arguments}
\label{sec:dedydadicization}

\subsection{The leading law follows by squeezing}

Let $x=2^{-t}$ and choose the integer $n$ with
$n\le t<n+1$.  Monotonicity gives
\begin{equation}
 F(2^{-(n+1)})\le F(x)\le F(2^{-n}).
\label{eq:monotone-squeeze}
\end{equation}
From \eqref{eq:integer-dyadic-expansion},
\begin{equation}
 \log F(2^{-n})=-\frac L2n^2+\bigO(n\log n).
\label{eq:dyadic-leading-bound}
\end{equation}
The same estimate holds with $n+1$, and $t=n+\bigO(1)$.  Dividing the logarithm of
\eqref{eq:monotone-squeeze} by $t^2$ gives
\begin{equation}
 \log F(2^{-t})\sim-\frac L2t^2.
\end{equation}
Since $\log x=-Lt$, this is \eqref{eq:global-quadratic-law}.  Thus the exact reciprocal-power sequence does recover the universal leading endpoint behavior for every real argument.

\subsection{Why squeezing cannot recover the sharp periodic formula}

The logarithmic gap between neighboring exact samples is of order $n$:
\begin{equation}
 \log F(2^{-n})-\log F(2^{-(n+1)})=Ln+\bigO(\log n).
\label{eq:neighbor-gap}
\end{equation}
It diverges.  Therefore monotonicity between consecutive reciprocal powers cannot determine even the constant term in $\log F(x)$, much less an error of order $1/\lambda$.

To approximate $x$ finely enough for an additive $o(1)$ error in $\log F(x)$, a dyadic grid needs relative spacing much smaller than $1/\lambda$.  At level $n$, that means a reduced numerator of size much larger than $\lambda$.  The fixed- or moderate-numerator analysis in \cref{sec:fixed-numerator} instead assumes that $r=\lambda/m$ is large.  The two regimes do not overlap.

\subsection{The remaining two-parameter problem}

For a general dyadic sequence $x_n=m_n2^{-n}$, the exact coefficient is
\begin{equation}
 F(x_n)=2^{-\binom n2}[z^n]\bigl(B_{m_n}(z)G(z)\bigr).
\end{equation}
When $m_n\gg\lambda_n$, the candidate radius $r_n=\lambda_n/m_n$ is small and the prefix polynomial
\begin{equation}
 \Prefix_{m_n}(r_n)=\sum_{h<m_n}\TM_h\e^{-hr_n}
\label{eq:large-prefix-problem}
\end{equation}
can no longer be replaced by $1$.  Its behavior depends on a long Thue--Morse prefix and, through binary refinement, on the numerator's digits.  A uniform saddle theorem in this regime would amount to a two-parameter de-Poissonization of the exact dyadic formula.  It is a legitimate route to the full continuous asymptotic, but it is no simpler than retaining the negative Laplace transform and applying the existing Bromwich analysis directly.

This is the precise boundary of the present re-derivation:
\begin{itemize}
\item exact coefficients recover the complete sharp expansion on reciprocal powers and all fixed-numerator rays;
\item exact coefficients plus monotonicity recover the leading quadratic law for all real $x$;
\item the full uniform sharp expansion for arbitrary $x$ requires control of growing prefix factors, which is the continuous problem in another coordinate system.
\end{itemize}

\section{A formalization roadmap}
\label{sec:formalization}

Most algebraic ingredients are already formalized in the repository.  The new bridge can be organized into a comparatively small analytic layer.

\begin{enumerate}
\item \textbf{Normalized coefficient sequence.}  Package
$a_n=d_n/n!=2^{\binom n2}F(2^{-n})$ and identify its ordinary generating function with $G$.

\item \textbf{Vertical coefficient extraction.}  Prove
\eqref{eq:vertical-extraction} from the bounded-support moment representation.  For $n\ge1$, absolute convergence follows from
$|G(c+\ii t)|\le G(c)$.

\item \textbf{Exact dyadic saddle mass.}  Define $\CoeffMass_n$ by
\eqref{eq:Cn-def} and prove the exact factorization
\eqref{eq:exact-coefficient-saddle}.

\item \textbf{Phase locking.}  From
$\lambda_n2^{-\lambda_n}=2^{-n}$ prove
$\log_2\lambda_n=\lambda_n-n$, and use periodicity to rewrite
$\Psi(\log_2\lambda_n)$ as $\Psi(\lambda_n)$.

\item \textbf{Exact separation.}  Formalize the algebraic identity
\eqref{eq:key-algebra-cancellation} and derive
\eqref{eq:exact-dyadic-separation}.

\item \textbf{Reuse the generic saddle machinery.}  Prove
\eqref{eq:coefficient-local-phase}.  Its main part is literally the phase already used by the all-orders endpoint modules; the only new tail is exponentially small in $\lambda$ rather than in $2^\lambda$.

\item \textbf{Prefix rays.}  Use the elementary bound
\eqref{eq:Sm-bound} to lift the same proof to fixed $m$, yielding
\eqref{eq:fixed-m-exact-separation} and \cref{thm:main-fixed-m}.

\item \textbf{Optional rescaled-Bromwich theorem.}  Formalize
\eqref{eq:negative-refinement-iterated} and the support collapse
\eqref{eq:positive-part-collapse}.  This provides an exact commutative square between the moment and endpoint inversion proofs.
\end{enumerate}

The conceptual gain is substantial: the asymptotic theorem at reciprocal powers would no longer be merely checked against exact values.  It would be derived from the exact-value infrastructure, with the continuous and discrete saddle phases proved equal.

\section{Conclusions}
\label{sec:conclusions}

The exact arithmetic and endpoint asymptotic of the Fabius function are connected more tightly than their usual presentations suggest.

The reciprocal-power formula first removes a universal factor
$2^{-\binom n2}$ and leaves a coefficient $[z^n]G(z)$.  The recurrence and the composition formula resum into the entire dyadic product for $G$.  Its negative-real restriction is exactly the Laplace product whose Mellin transform produces the quadratic drift, the constant, and the Gamma--zeta periodic term.  Coefficient extraction at the Lambert radius then supplies the Gaussian normalization and the inverse-power corrections.

The crucial arithmetic identity
\[
 n=\lambda_n-\log_2\lambda_n
\]
has two simultaneous consequences.  It makes $\lambda_n$ the coefficient saddle, and it locks the log-periodic phase:
$\Psi(\log_2\lambda_n)=\Psi(\lambda_n)$.  Once this is noticed, the local coefficient phase becomes identical to the local continuous Bromwich phase.  The equality persists to every order.

The Gaussian-binomial formula contributes a complementary lesson.  Its conspicuous
$q$-Pochhammer normalizer is not itself the endpoint asymptotic.  Exact geometric interpolation cancels it against the Mersenne factor in the numerator and exposes the coefficient
$[z^n](B_mG)$.  For fixed $m$, the prefix factor has one dominant exponential, and the same saddle analysis gives the full sharp expansion on $m2^{-n}$.

Finally, rescaling the continuous inversion at $2^{-n}$ and iterating the negative refinement equation maps the saddle $2^{\lambda_n}$ to $\lambda_n$ exactly.  The finite Thue--Morse product introduced by those refinement steps is exponentially close to $1$ at the saddle and disappears exactly after integration because the probability law is supported on $[0,1]$.  This identity supplies a direct analytic explanation of why the discrete and continuous derivations agree.

Thus the missing connection can be stated succinctly:
\begin{equation*}
 \boxed{\begin{gathered}
 \text{finite dyadic arithmetic}\\
 \Big\downarrow\ \text{ exact cancellation and resummation}\\
 \text{the same infinite dyadic product}\\
 \Big\downarrow\ \text{ coefficient saddle}\\
 \text{the corrected endpoint asymptotic}
 \end{gathered}}
\end{equation*}

\appendix

\section{The key algebraic cancellation}
\label{app:algebra}

For completeness, let $p=L(\lambda-n)$.  Then $n=\lambda-p/L$, and
\begin{align*}
 -\frac L2(n^2-n)
 &=-\frac L2\lambda^2+\lambda p-\frac{p^2}{2L}
 +\frac L2\lambda-\frac p2,\\
 -np&=-\lambda p+\frac{p^2}{L}.
\end{align*}
Therefore
\begin{align*}
 &-\frac L2(n^2-n)+\lambda-\frac{p^2}{2L}+\frac p2-np\\
 &\quad=-\frac L2\lambda^2+
 \left(1+\frac L2\right)\lambda.
\end{align*}
For reciprocal powers, $p=\log\lambda$.  For a fixed numerator $m$, the same calculation uses
$p=\log(\lambda/m)$, because
$\log_2(\lambda/m)=\lambda-n$.

\section{Exact arithmetic used in the numerical tables}
\label{app:numerics}

The reciprocal coefficients are generated by
\begin{equation}
 a_0=1,
 \qquad
 a_n=\frac1{2^n-1}
 \sum_{j=0}^{n-1}\frac{a_j}{(n-j+1)!}.
\end{equation}
All $a_n$ are stored as reduced rational numbers, and
\begin{equation}
 F(2^{-n})=2^{-\binom n2}a_n.
\end{equation}
For a fixed numerator,
\begin{align}
 [z^n](B_m(z)G(z))
 &=\sum_{h=0}^{m-1}\TM_h
 \sum_{j=0}^{n}a_j
 \frac{(m-1-h)^{n-j}}{(n-j)!},                             \label{eq:exact-prefix-coeff-compute}\\
 F(m2^{-n})
 &=2^{-\binom n2}[z^n](B_mG).                              \label{eq:exact-prefix-value-compute}
\end{align}
Equations \eqref{eq:exact-prefix-coeff-compute}--\eqref{eq:exact-prefix-value-compute} involve rational arithmetic only.  The Lambert coordinate is then evaluated at high precision from
\eqref{eq:lambda-def}; $\Psi$ and its derivatives are evaluated from
\eqref{eq:Psi-Fourier}.  Because the Fourier coefficients decay exponentially, a small number of modes suffices for the digits printed in \cref{tab:inverse-numerics,tab:fixed-m-numerics}.

\section{Bibliographic note}

The exact arithmetic identities used here are developed in the ProveIt repository and its accompanying exposition \cite{proveit,masterdoc}, building in particular on the compact-support and arithmetic papers of Arias de Reyna \cite{arias1982,arias2018}.  Haugland \cite{haugland2016} gives another exact-evaluation perspective.  The Mellin treatment of harmonic sums follows the general framework of Flajolet, Gourdon, and Dumas \cite{flajolet1995}, while the Lambert coordinate uses the branch conventions of Corless et al. \cite{corless1996}.  The contribution of this report is the explicit coefficient-saddle bridge, its phase-locking identity, its extension to fixed dyadic numerators, and the exact rescaled-Bromwich collapse.

\begin{thebibliography}{99}

\bibitem{proveit}
V.~Reshetnikov,
\newblock \emph{ProveIt: Analysis/FabiusFunction}, Lean formal development,
\newblock GitHub repository, snapshot accessed 25 August 2026,
\newblock \url{https://github.com/VladimirReshetnikov/ProveIt/tree/main/Analysis/FabiusFunction}.

\bibitem{masterdoc}
V.~Reshetnikov,
\newblock \emph{The Fabius Function and Rvachev's Up-Function: Exact arithmetic, probability, Fourier analysis, and corrected sharp asymptotics},
\newblock repository exposition, snapshot accessed 25 August 2026,
\newblock \url{https://github.com/VladimirReshetnikov/ProveIt/blob/main/Analysis/FabiusFunction/docs/Fabius_Function_and_Rvachev_Up/Fabius_Function_and_Rvachev_Up.tex}.

\bibitem{arias1982}
J.~Arias de Reyna,
\newblock Definici\'on y estudio de una funci\'on indefinidamente diferenciable de soporte compacto,
\newblock \emph{Revista de la Real Academia de Ciencias Exactas, F\'isicas y Naturales de Madrid}
\textbf{76} (1982), no.~1, 21--38.
\newblock English translation: \emph{An infinitely differentiable function with compact support: Definition and properties}, arXiv:1702.05442 (2017).

\bibitem{arias2018}
J.~Arias de Reyna,
\newblock Arithmetic of the Fabius function,
\newblock \emph{INTEGERS} \textbf{18} (2018), Article A51; arXiv:1702.06487, version~3.

\bibitem{haugland2016}
J.~K.~Haugland,
\newblock Evaluating the Fabius function,
\newblock arXiv:1609.07999 (2016).

\bibitem{flajolet1995}
P.~Flajolet, X.~Gourdon, and P.~Dumas,
\newblock Mellin transforms and asymptotics: harmonic sums,
\newblock \emph{Theoretical Computer Science} \textbf{144} (1995), 3--58.

\bibitem{corless1996}
R.~M.~Corless, G.~H.~Gonnet, D.~E.~G.~Hare, D.~J.~Jeffrey, and D.~E.~Knuth,
\newblock On the Lambert $W$ function,
\newblock \emph{Advances in Computational Mathematics} \textbf{5} (1996), 329--359.

\end{thebibliography}

\end{document}
